\documentclass[12pt,a4paper]{report}

\usepackage[T1]{fontenc}
\usepackage[utf8]{inputenc}
\usepackage[brazil]{babel}
\usepackage{geometry}
\usepackage{setspace}
\usepackage{microtype}
\usepackage{amsmath,amssymb}
\usepackage{booktabs}
\usepackage{longtable}
\usepackage{array}
\usepackage{siunitx}
\usepackage{graphicx}
\usepackage{hyperref}
\usepackage[numbers,sort&compress]{natbib}
\usepackage{url}

\geometry{margin=2.5cm}
\setstretch{1.15}
\setlength{\parindent}{1.25cm}
\setlength{\parskip}{0.2em}
\sisetup{locale = DE}
\urlstyle{same}
\hypersetup{
  colorlinks=true,
  linkcolor=black,
  citecolor=black,
  urlcolor=blue,
  pdftitle={CabotageLens: avaliacao multimodal de custos, combustivel e emissoes no transporte de cargas no Brasil},
  pdfauthor={Projeto CabotageLens}
}

\begin{document}

\begin{titlepage}
\centering
{\Large\bfseries CabotageLens: avaliação multimodal de custos, combustível e emissões no transporte de cargas no Brasil\par}
\vspace{1.2cm}
{\large Relatório técnico-científico em formato de trabalho final, redigido a partir do repositório \texttt{cabotage-lens}\par}
\vspace{0.8cm}
{\large Formulação metodológica, base empírica e implementação computacional\par}
\vfill
{\large \today\par}
\end{titlepage}

\pagenumbering{roman}

\chapter*{Resumo}
\addcontentsline{toc}{chapter}{Resumo}
Este trabalho documenta, em formato de relatório técnico-científico, a formulação metodológica e a implementação computacional do projeto CabotageLens, voltado à comparação \emph{door-to-door} entre transporte rodoviário e alternativas multimodais com cabotagem no Brasil. O problema de pesquisa parte da elevada dependência da matriz brasileira de cargas em relação ao modo rodoviário e da necessidade de um instrumento reprodutível para estimar, com rastreabilidade, combustível, emissões e custo energético em corredores candidatos à substituição modal. A principal contribuição metodológica do trabalho é a construção de uma base observada de cabotagens conteinerizadas brasileiras a partir da ANTAQ, segmentada por viagem e por parada portuária, e sua combinação com indicadores do EU MRV associados por IMO, com o objetivo de estimar uma mediana de intensidade por par origem--destino ponderada pelo trabalho de transporte de todas as viagens elegíveis. No recorte operacional de 2025, a base da ANTAQ reúne \num{1324} viagens, \num{6797} paradas, \num{7103} ocorrências de chamadas portuárias e \num{389} IMOs únicos. Desses IMOs, \num{243} foram associados com sucesso a indicadores positivos de consumo por \emph{transport work} na base MRV, cobrindo \num{62.5}\% dos IMOs, \num{59.5}\% das viagens, \num{52.9}\% da carga observada em toneladas e \num{50.6}\% da carga observada em TEU. Para as demais viagens, o cálculo registra uma estatística de classe ou tipo de navio como fallback, com regra explícita de aparagem de caudas e proveniência da amostra. A matriz marítima contém \num{177} pares direcionais, construídos a partir de corredores completos observados dentro da mesma viagem e do somatório de seus subtrechos com carga a bordo. O texto também descreve a unidade operacional de \num{1} contêiner e \num{14} toneladas e o uso programático da planilha de referência associada aos estudos de Gustavo Costa. A formulação resultante é híbrida e auditável, com limitações relacionadas à cobertura exata por IMO, à observabilidade de rotas e ao escopo predominantemente \emph{tank-to-wheel} do modelo.

\noindent\textbf{Palavras-chave:} cabotagem; transporte multimodal; ANTAQ; EU MRV; emissões; combustível; logística.

\chapter*{Abstract}
\addcontentsline{toc}{chapter}{Abstract}
This report documents the methodological design and software implementation of CabotageLens, a computational tool built to compare all-road freight transport with multimodal alternatives that include Brazilian coastal shipping. The research problem stems from the strong dependence of Brazilian freight flows on road transport and from the lack of reproducible tools able to estimate fuel use, direct emissions, and energy-related cost in a transparent door-to-door manner. The core methodological contribution is the combination of an observed 2025 cabotage dataset built from ANTAQ port-call and cargo records with ship-level EU MRV indicators matched by IMO, in order to estimate an origin--destination maritime fuel-intensity median weighted by transport work across all eligible voyages. The ANTAQ-derived base contains \num{1324} voyages, \num{6797} stops, \num{7103} port-call occurrences, and \num{389} unique IMOs. Among those IMOs, \num{243} were successfully matched to positive MRV transport-work fuel indicators, covering \num{62.5}\% of unique IMOs, \num{59.5}\% of voyages, \num{52.9}\% of observed loaded weight, and \num{50.6}\% of observed loaded TEU. The remaining voyages use a recorded vessel-class or ship-type fallback with explicit tail trimming and sample provenance. The maritime matrix contains \num{177} directional port pairs built from complete corridors observed within the same voyage and the sum of their cargo-aware sublegs. The report also defines a benchmark cargo of one container and fourteen tonnes and a programmatic comparison with the workbook associated with Gustavo Costa's studies. The resulting pipeline is a hybrid, auditable formulation subject to exact-IMO coverage limits, route observability constraints, and a predominantly tank-to-wheel scope.

\noindent\textbf{Keywords:} coastal shipping; multimodal freight; ANTAQ; EU MRV; emissions; fuel use; logistics.

\tableofcontents
\clearpage
\pagenumbering{arabic}

\chapter{Introdução}
\label{chap:introducao}

O transporte de cargas no Brasil é fortemente marcado pela dominância do modo rodoviário. Essa configuração produz um problema operacional e analítico relevante: para muitos pares origem--destino de média e longa distância, o modo predominante pode não ser o mais eficiente do ponto de vista energético ou climático. Ao mesmo tempo, a simples afirmação de que ``a cabotagem emite menos'' é insuficiente, porque a comparação correta não é entre um navio isolado e um caminhão isolado, e sim entre cadeias logísticas completas, com acessos terrestres, operações portuárias, restrições geográficas e diferentes níveis de utilização de capacidade \citep{icct2022,shortsea2019,competitiveness2024}.

No contexto brasileiro, esse problema é ainda mais sensível. A cabotagem conteinerizada depende de uma combinação específica de oferta de serviço, conectividade terrestre, localização portuária e custos de interface. O estudo de competitividade em super-rede dedicado à cabotagem brasileira mostra que a atratividade do modo marítimo emerge de uma lógica de rede multimodal e não de uma vantagem absoluta e uniforme sobre o transporte rodoviário \citep{competitiveness2024}. De forma complementar, o estudo de trajetórias de descarbonização da cabotagem brasileira reforça que a análise do desempenho ambiental do setor depende da unidade funcional escolhida, do combustível considerado e da forma como o consumo do navio é atribuído à carga analisada \citep{decarb2024}. O problema de pesquisa deste trabalho, portanto, pode ser formulado nos seguintes termos: como comparar, de modo reprodutível e auditável, o desempenho de uma rota rodoviária direta com o de uma alternativa rodoviária--cabotagem--rodoviária para cargas conteinerizadas no Brasil?

O projeto CabotageLens responde a essa questão por meio de uma cadeia completa de avaliação que resolve coordenadas de origem e destino, calcula rotas rodoviárias, escolhe portos de acesso, consulta ou reutiliza distâncias marítimas, estima combustível e emissões por perna e consolida os resultados em artefatos persistentes e interfaces de análise. O modelo incorpora dados observados de cabotagem brasileira e indicadores empíricos do EU MRV associados por IMO, permitindo construir métricas marítimas por rota e utilizar fatores de classe ou tipo apenas como fallback identificado.

O objetivo geral deste relatório é documentar a formulação metodológica e computacional com a forma e o rigor de um trabalho final acadêmico. Os objetivos específicos são quatro. Primeiro, situar teoricamente o problema de comparação modal entre rodoviário e cabotagem à luz da literatura local e internacional. Segundo, documentar a base empírica usada no projeto, com destaque para a reconstrução de viagens da ANTAQ e para a extração de eficiência de combustível do EU MRV. Terceiro, explicitar como esses dados são transformados em indicadores operacionais e em avaliações \emph{door-to-door}. Quarto, apresentar a cobertura do cruzamento ANTAQ+MRV, a \texttt{sea\_matrix} direcional e o benchmark com a planilha de referência enviada por Gustavo Costa.

O trabalho também se justifica por um motivo metodológico mais amplo. Grande parte das comparações modais disponíveis em relatórios técnicos utiliza fatores médios fixos por tonelada-quilômetro, sem explicitar suficientemente de onde vieram esses fatores, a que frota se referem ou como foram adaptados à remessa observada. O CabotageLens adota uma abordagem auditável: a unidade de cálculo é a perna logística; a unidade observacional marítima é o trecho portuário operado; e a intensidade representativa é a mediana das intensidades resolvidas para todas as viagens que ligaram o par observado pela ANTAQ, ponderada pelo trabalho de transporte de cada viagem.

Este documento está organizado da seguinte maneira. O Capítulo~\ref{chap:referencial} apresenta o referencial teórico e os trabalhos correlatos. O Capítulo~\ref{chap:base} descreve a base de dados, os artefatos processados e a organização do repositório. O Capítulo~\ref{chap:metodo} formaliza a metodologia adotada. O Capítulo~\ref{chap:resultados} discute as principais evidências empíricas. O Capítulo~\ref{chap:implementacao} sintetiza a implementação computacional. O Capítulo~\ref{chap:limitacoes} discute limitações e ameaças à validade. Por fim, o Capítulo~\ref{chap:conclusao} apresenta as conclusões.

\chapter{Referencial teórico e trabalhos correlatos}
\label{chap:referencial}

\section{Cabotagem, mudança modal e competitividade logística}

O debate sobre cabotagem brasileira costuma oscilar entre duas narrativas. A primeira é a narrativa de potencial: em um país com grande litoral, corredores longos e concentração urbana na faixa costeira, a cabotagem poderia absorver parte dos fluxos de carga hoje realizados por rodovia. A segunda é a narrativa de fricção: a mesma cabotagem enfrenta barreiras de integração modal, frequência de serviço, custos portuários e incertezas operacionais que impedem uma transição modal automática \citep{icct2022,users2022}. O interesse científico do problema está justamente nessa tensão.

O artigo sobre competitividade da cabotagem brasileira em uma super-rede multimodal é particularmente importante para este trabalho porque mostra que a questão correta não é se o modo marítimo é ``melhor'' em abstrato, mas em quais relações origem--destino, com que estrutura de acesso e em que topologia de rede ele se torna competitivo \citep{competitiveness2024}. Esse enquadramento é totalmente compatível com o CabotageLens: o projeto não procura atribuir um único fator médio a toda a cabotagem, mas avaliar rotas, pernas e combinações concretas de origem, porto e destino.

\section{Eficiência energética marítima e uso de dados MRV}

A literatura baseada no MRV europeu oferece indicadores reportados de consumo e emissões em operação, reduzindo a dependência exclusiva de coeficientes teóricos ou valores de projeto \citep{eumrv2025}. Estudos de monitoramento da pegada de carbono com dados MRV demonstram a utilidade desses registros para avaliar desempenho observado, inclusive em outros segmentos de frota \citep{drybulk2019}. Para o transporte conteinerizado, a intensidade adequada à comparação com a carga não é apenas combustível por distância, mas combustível por \emph{transport work}, isto é, por tonelada transportada e por milha náutica.

Este trabalho adota essa lógica. O indicador central extraído do MRV não é o consumo bruto anual do navio, mas o consumo médio anual por \emph{transport work}. Essa escolha tem duas implicações analíticas importantes. Primeiro, ela embute uma normalização pelo trabalho efetivamente realizado pelo navio. Segundo, ela permite atribuir combustível à carga transportada sem recorrer, como regra principal, a suposições arbitrárias sobre participação percentual da remessa no consumo total da embarcação.

\section{Operações portuárias, hoteling e escopo ambiental}

Embora o trecho marítimo seja o centro da comparação modal, a literatura deixa claro que uma avaliação \emph{door-to-door} não pode ignorar as interfaces portuárias e o consumo fora da navegação de cruzeiro. Estudos de emissões em berço e durante carregamento e descarregamento mostram que fases de hoteling e operação podem representar parcelas relevantes do inventário, sobretudo quando a permanência em porto é elevada \citep{berth2009,shipops2022}. Em paralelo, trabalhos sobre consumo de RTGs e sistemas híbridos de pátio ajudam a justificar a modelagem por equipamento nos terminais \citep{rtg2017,hybridrtg2021}.

No CabotageLens, esses elementos permanecem metodologicamente relevantes, mas sua utilização depende do tipo de KPI marítimo empregado. Quando a avaliação usa o indicador direcional derivado do MRV por \emph{transport work}, o sistema adota a hipótese de fronteira de que a intensidade anual já abrange o consumo reportado do navio e evita somar hoteling novamente para reduzir o risco de dupla contagem. Já nas avaliações de contingência, sensibilidade ou decomposição por classe de navio, os submodelos de hoteling e de operações portuárias continuam desempenhando papel importante. Essa distinção é coerente com a literatura de inventário e com o próprio guia metodológico EMEP/EEA, citado internamente no projeto como base para a derivação das taxas auxiliares de hoteling \citep{emep2023}.

\section{Escopo TTW, combustíveis e limites de interpretação}

O projeto está metodologicamente mais próximo de um modelo \emph{tank-to-wheel} do que de uma avaliação completa de ciclo de vida. Essa escolha é deliberada. A literatura recente sobre combustíveis marítimos deixa claro que conclusões \emph{well-to-wake} ou \emph{life cycle assessment} exigem fatores adicionais, hipóteses de suprimento e tratamento explícito de incertezas \citep{lifecycle2020,maritimelca2024}. O trabalho de descarbonização da cabotagem brasileira também reforça que a comparação entre VLSFO, HVO, LNG e biodiesel depende da fronteira analítica escolhida \citep{decarb2024,biodiesel2017}.

Assim, este relatório adota a seguinte posição: o CabotageLens é uma ferramenta para comparação operacional de combustível, custo energético e emissões diretas de combustão. Ele não substitui estudos completos de ciclo de vida, mas fornece uma base auditável para análises comparativas e para cenários de combustíveis alternativos.

\chapter{Base empírica, dados e artefatos do projeto}
\label{chap:base}

\section{Organização do repositório}

O repositório está organizado de forma a separar claramente lógica de negócio, interfaces, persistência e artefatos de dados. O diretório \texttt{modules/} concentra os componentes principais de roteamento, avaliação rodoviária, avaliação marítima, custos, emissões, persistência e serviços auxiliares. O diretório \texttt{app/} expõe a aplicação Streamlit, enquanto \texttt{scripts/} oferece fluxos reproduzíveis de linha de comando para avaliação unitária, avaliação em lote, sincronização com Supabase Storage e materialização de tabelas derivadas. O diretório \texttt{data/} armazena insumos versionados e artefatos processados não transacionais. Já o diretório \texttt{supabase/} registra as migrações SQL correspondentes às tabelas persistidas do projeto.

Essa organização tem implicação metodológica direta: a geração de resultados não depende de manipulações ad hoc externas ao código versionado. Os artefatos centrais podem ser refeitos a partir de scripts determinísticos, e os resultados operacionais podem ser reconsultados via banco ou via arquivos persistidos em Storage. Em outras palavras, a estrutura do repositório reforça a auditabilidade do modelo.

\section{Fontes primárias e materiais de apoio}

Quatro grupos de fontes são centrais para o projeto:

\begin{longtable}{p{3.3cm} p{4.4cm} p{7.3cm}}
\caption{Principais fontes de dados e seu papel metodológico}
\label{tab:fontes-dados}\\
\toprule
Fonte & Artefato principal & Papel no modelo \\
\midrule
\endfirsthead
\toprule
Fonte & Artefato principal & Papel no modelo \\
\midrule
\endhead
ANTAQ & \texttt{2025Atracacao.txt}, \texttt{2025Carga.txt}, \texttt{2025TemposAtracacao.txt} & Reconstrução observada de viagens, paradas, sequência portuária, carga movimentada e tempos de estadia por escala \citep{antaq2025}. \\
EU MRV & Workbooks anuais de publicação de informação (\texttt{2018}--\texttt{2024}) & Extração de indicadores empíricos de eficiência marítima, com foco em combustível por \emph{transport work} e identificação dos navios por IMO \citep{eumrv2025,drybulk2019}. \\
ANP & \texttt{semanal-estados-desde-2013.xlsx} e \texttt{latest\_diesel\_prices.csv} & Parametrização do custo do diesel rodoviário por unidade da federação \citep{anp2026}. \\
Workbook de referência & \texttt{docs/references/core/Dados Relatorio 2.xlsx} & Benchmark operacional local-only e apoio à modelagem de cenários de cabotagem, operações portuárias e comparação com resultados de referência associados aos estudos de Gustavo Costa \citep{workbookdados,decarb2024}. \\
\bottomrule
\end{longtable}

Além dessas fontes, o projeto utiliza \texttt{ports\_br.json} como cadastro portuário nacional estruturado e \texttt{sea\_matrix.json} como matriz de distâncias marítimas entre portos brasileiros. Esses dois artefatos funcionam como camada geométrica e topológica básica sobre a qual a avaliação multimodal é construída.

\section{Artefatos processados centrais}

Os artefatos processados centrais para a lógica marítima são:

\begin{longtable}{p{5.2cm} p{9.0cm}}
\caption{Artefatos processados centrais do CabotageLens}
\label{tab:artefatos-processados}\\
\toprule
Artefato & Função no projeto \\
\midrule
\endfirsthead
\toprule
Artefato & Função no projeto \\
\midrule
\endhead
\texttt{container\_ship\_efficiency\_classes.json} & Consolida classes de navios porta-contêineres e suas estatísticas robustas de consumo por distância e por \emph{transport work}. \\
\texttt{container\_ship\_fuel\_rate\_sea\_by\_class.json} & Mantém taxas de consumo marítimo por classe de navio para usos de contingência e decomposição por classe. \\
\texttt{container\_ship\_hoteling\_rate\_by\_class.json} & Armazena taxas horárias de hoteling derivadas do guia EMEP/EEA \citep{emep2023}. \\
\texttt{antaq\_voyages.csv} & Tabela normalizada com uma linha por viagem observada na ANTAQ. \\
\texttt{antaq\_voyage\_stops.csv} & Tabela normalizada com uma linha por parada observada, incluindo carga e tempos portuários. \\
\texttt{antaq\_voyage\_stop\_calls.csv} & Tabela normalizada com granularidade de \texttt{call\_id}. \\
\texttt{mrv\_average\_efficiency\_by\_imo.json} & Lookup por IMO dos indicadores MRV aproveitáveis para os navios observados na cabotagem. \\
\texttt{sea\_matrix.json} & Matriz porto--porto de distâncias marítimas enriquecida com mediana de intensidade ponderada por trabalho de transporte para cada par origem--destino e corredores observados para caminho e distância. \\
\bottomrule
\end{longtable}

Na prática, esses artefatos podem residir localmente em \texttt{data/processed/...} ou ser resolvidos via Supabase Storage, mas o conteúdo lógico é o mesmo. O projeto foi deliberadamente desenhado para desacoplar cálculo, persistência e interface sem perder reprodutibilidade.

\section{Base observada de cabotagens conteinerizadas em 2025}

A base observada de viagens de cabotagem conteinerizada é construída a partir da ANTAQ. O pipeline cruza registros de atracação, carga e tempos de atracação, colapsa chamadas consecutivas no mesmo porto, segmenta as viagens por IMO e consolida cada viagem em uma sequência ordenada de origem, paradas intermediárias e destino.

No recorte de 2025, a base processada contém:

\begin{table}[htbp]
\centering
\caption{Escala da base observada de cabotagens conteinerizadas}
\label{tab:base-observada}
\begin{tabular}{lr}
\toprule
Indicador & Valor \\
\midrule
Viagens observadas & \num{1324} \\
Paradas observadas & \num{6797} \\
Ocorrências de chamadas portuárias & \num{7103} \\
IMOs únicos & \num{389} \\
Carga carregada total (t) & \num{30191845.948} \\
Carga carregada total (TEU) & \num{2872715.0} \\
\bottomrule
\end{tabular}
\end{table}

Esses números caracterizam a base observada de operação e fornecem a matéria-prima para uma lógica baseada em navios, viagens e trechos efetivamente registrados. As estatísticas de classe permanecem disponíveis como fallback identificado quando não existe associação positiva por IMO.

\chapter{Método}
\label{chap:metodo}

\section{Unidade funcional e cenários comparados}

O CabotageLens compara dois cenários logísticos para um mesmo par origem--destino. O primeiro é o cenário \emph{road-only}, no qual toda a remessa é transportada por rodovia. O segundo é o cenário multimodal, composto por \emph{first mile} rodoviário até o porto de origem, perna marítima de cabotagem entre portos brasileiros e \emph{last mile} rodoviário até o destino final.

Como benchmark operacional padrão do projeto, a avaliação unitária adota \num{1} contêiner e \num{14} toneladas de carga. Essa escolha aproxima a ferramenta do benchmark operacional usado nos materiais de referência ligados aos estudos de Gustavo Costa e evita a ambiguidade de comparar cenários com remessas implícitas de massa excessivamente alta para uma unidade conteinerizada simples.

\section{Modelo rodoviário}

O modelo rodoviário parte da massa da carga, do \emph{preset} de veículo e da distância roteirizada por OpenRouteService. Seja \(M\) a carga total em toneladas e \(P\) a capacidade útil do caminhão selecionado. O número de viagens necessárias é dado por:
\begin{equation}
n_{\text{viagens}} = \left\lceil \frac{M}{P} \right\rceil.
\label{eq:road-trips}
\end{equation}

Se \(d_{\text{rod}}\) representa a distância do trecho em quilômetros e \(\eta_{\text{rod}}\) a eficiência ajustada do veículo em km/L, o consumo de diesel do trecho rodoviário é:
\begin{equation}
L_{\text{rod}} = \frac{d_{\text{rod}}}{\eta_{\text{rod}}} \cdot n_{\text{viagens}}.
\label{eq:road-fuel}
\end{equation}

Com \(p_{\text{diesel}}\) como preço de diesel em BRL/L, obtido a partir das referências por unidade da federação, o custo energético e as emissões diretas são calculados por:
\begin{equation}
C_{\text{rod}} = L_{\text{rod}} \cdot p_{\text{diesel}}
\qquad\text{e}\qquad
E_{\text{rod}} = L_{\text{rod}} \cdot EF_{\text{diesel}},
\label{eq:road-cost-emissions}
\end{equation}
em que \(EF_{\text{diesel}}\) é o fator de emissão direta do diesel adotado no projeto.

Essa formulação não esgota a complexidade do transporte rodoviário, mas oferece uma base transparente, reexecutável e compatível com a lógica do repositório: toda perna de cálculo precisa ser rastreável até parâmetros explícitos e dados de entrada conhecidos.

\section{Reconstrução observada de viagens da ANTAQ}

O passo metodológico mais importante do lado marítimo é a reconstrução das viagens observadas. O pipeline associa registros de carga a chamadas portuárias, filtra o recorte de cabotagem conteinerizada, exige identificação de IMO, colapsa chamadas consecutivas no mesmo porto e segmenta a trajetória por navio. Como resultado, cada viagem observada possui:

\begin{enumerate}
  \item um identificador único de viagem;
  \item um IMO;
  \item uma sequência ordenada de paradas;
  \item carga carregada, descarregada, movimentada e líquida por parada;
  \item tempos observados de permanência, operação e espera em porto.
\end{enumerate}

Essa estrutura é relevante porque permite responder à pergunta operacional correta para o problema da cabotagem: qual navio, em qual trajetória, carregou qual quantidade de carga entre quais pares consecutivos de portos?

\section{Extração de eficiência marítima por IMO a partir do EU MRV}

O projeto processa os workbooks anuais do EU MRV e extrai, para navios do tipo \emph{container ship}, o indicador de consumo médio anual por \emph{transport work}. Na prática, a prioridade é dada ao campo de combustível por trabalho de transporte baseado em massa carregada, pois ele já representa combustível por tonelada transportada e por milha náutica. Esse KPI é superior, para a finalidade deste trabalho, ao uso de consumo absoluto anual do navio, porque já normaliza o consumo pelo serviço efetivamente prestado.

Depois da extração, os navios observados na ANTAQ são associados por IMO a esses indicadores MRV. O resultado consolidado do cruzamento é resumido na Tabela~\ref{tab:cobertura-mrv}.

\begin{table}[htbp]
\centering
\caption{Cobertura do cruzamento ANTAQ + EU MRV}
\label{tab:cobertura-mrv}
\begin{tabular}{lrr}
\toprule
Indicador & Valor absoluto & Cobertura \\
\midrule
IMOs com match MRV & \num{243} de \num{389} & \num{62.5}\% \\
Viagens com match MRV & \num{788} de \num{1324} & \num{59.5}\% \\
Carga carregada com match MRV (t) & \num{15959761.561} de \num{30191845.948} & \num{52.9}\% \\
Carga carregada com match MRV (TEU) & \num{1454351.75} de \num{2872715.0} & \num{50.6}\% \\
\bottomrule
\end{tabular}
\end{table}

Esses números mostram que a cobertura exata por IMO é parcial. Os estimadores direcionais por par preservam essa condição ao identificar separadamente as viagens sem correspondência como fallback de classe ou tipo.

\section{Construção do KPI marítimo direcional por rota}

A existência de um KPI por navio não resolve isoladamente o problema do modelo, porque o objetivo da ferramenta é avaliar pares porto--porto ou origem--destino. Para isso, o projeto percorre a sequência normalizada de paradas de cada viagem e retém o corredor completo sempre que a origem aparece antes do destino. Cada \texttt{voyage\_id} contribui uma vez para cada par direcional; se chamadas repetidas produzirem mais de um recorte elegível, aplica-se dentro da viagem o mesmo critério do modelo: corredor direto primeiro e, na ausência dele, menor distância completa.

Considere uma viagem \(v\) e o conjunto \(\mathcal{S}_{v,o,d}\) de subtrechos contíguos entre a origem \(o\) e o destino \(d\). Para cada subtrecho \(s\), sejam \(d_s\) a distância marítima em milhas náuticas, \(m_s\) a carga reconstruída a bordo e \(I_v\) a intensidade resolvida para o navio, em g/(t$\cdot$nm). O trabalho de transporte da viagem no corredor é:
\begin{equation}
W_{v,o,d} = \sum_{s \in \mathcal{S}_{v,o,d}} m_s \cdot d_s.
\label{eq:segment-weight}
\end{equation}

O consumo derivado para a atividade conteinerizada observada nessa viagem é calculado pela soma dos subtrechos:
\begin{equation}
F_{v,o,d} = \frac{1}{1000}
\sum_{s \in \mathcal{S}_{v,o,d}} I_v \cdot m_s \cdot d_s.
\label{eq:route-intensity}
\end{equation}
Viagens com a mesma sequência portuária formam uma alternativa de corredor. Sua intensidade diagnóstica é a razão entre o combustível somado e o trabalho de transporte somado. A intensidade aplicada ao cenário, entretanto, usa todas as observações viagem--OD resolvidas, independentemente do corredor. Ordenando suas intensidades \(I_v\) e usando \(W_{v,o,d}\) como peso, define-se a mediana ponderada:
\begin{equation}
I_{o,d}^{\mathrm{rep}} = \min\left\{x:\sum_{v:I_v\leq x} W_{v,o,d} \geq \frac{1}{2}\sum_v W_{v,o,d}\right\}.
\label{eq:od-weighted-median}
\end{equation}
Observações com trabalho zero não influenciam a mediana quando existe trabalho positivo. Se todas as viagens resolvidas do par tiverem trabalho zero, usa-se uma mediana simples explicitamente marcada. Essa convenção evita a instabilidade da média aritmética diante de caudas elevadas e mantém o estimador independente de um corredor de referência.

Na resolução da intensidade, a primeira fonte é o IMO no EU MRV. Sem correspondência positiva, utiliza-se a média aparada em 1\% por cauda registrada para a classe, com mediana apenas quando esse campo não existe. Para o tipo de navio, ordenam-se os últimos valores positivos por IMO, exclui-se de cada cauda \(\lfloor 0{,}01n \rfloor\) observações e calcula-se a média da amostra retida; quando a amostra é pequena demais para remover ao menos uma observação por cauda, usa-se a mediana. Para o recorte conteinerizado sem metadado de tipo, \emph{container ship} é o fallback documentado. Valores associados diretamente ao IMO nunca são aparados. A proveniência por tipo registra fonte, estimador, regra, tamanhos bruto e retido, contagem excluída, limites e resumos não aparados; para classe, registram-se a estatística e os metadados de amostra disponíveis no artefato correspondente. O modelo avalia todas as sequências observadas dentro de uma mesma viagem, sem montar um caminho sintético com pernas de viagens diferentes. A seleção de um corredor direto utilizável ou, na sua ausência, do corredor completo de menor distância define somente o caminho e a distância; a intensidade representativa continua baseada em todas as viagens elegíveis do par.

O enriquecimento direcional da \texttt{sea\_matrix} resulta em:

\begin{table}[htbp]
\centering
\caption{Síntese do enriquecimento direcional da \texttt{sea\_matrix}}
\label{tab:sea-matrix-summary}
\begin{tabular}{lr}
\toprule
Indicador & Valor \\
\midrule
Pares direcionais com KPI & \num{177} \\
Nós de origem com ao menos um par válido & \num{16} \\
Trechos candidatos antes do colapso de aliases & \num{5473} \\
Trechos canônicos calculados & \num{5438} \\
Trechos com distância da matriz & \num{5023} \\
Trechos com fallback haversine & \num{415} \\
Trechos com intensidade por IMO & \num{3290} \\
Trechos com fallback de intensidade & \num{2148} \\
Observações viagem--OD preservadas & \num{12025} \\
Taxa de resolução de intensidade & \num{100.0}\% \\
\bottomrule
\end{tabular}
\end{table}

O filtro é restrito a corredores presentes na sequência normalizada de uma mesma viagem. Nenhum subtrecho é descartado por nome de terminal, distância ausente ou intensidade não resolvida no artefato. Quatro nomes de terminais são associados a seus portos canônicos; chamadas consecutivas no mesmo complexo são colapsadas somente depois de somar seus movimentos de carga. Para pares canônicos distintos sem distância positiva na matriz, o fallback haversine é registrado como aproximação. As sequências normalizadas derivam das chamadas com movimentação conteinerizada disponíveis no pipeline e, portanto, não devem ser interpretadas como inventário de toda atracação física sem carga observada.

\section{Lógica de avaliação marítima no motor do projeto}

A avaliação unitária carrega da \texttt{sea\_matrix} enriquecida o corredor completo selecionado no processamento offline. Para uma carga de cenário \(M\), o combustível marítimo de navegação é calculado pela soma das distâncias de seus subtrechos e pela intensidade representativa do par:
\begin{equation}
F_{\text{mar}} = \frac{1}{1000}
\sum_{s \in \mathcal{S}^{*}_{o,d}} I_{o,d}^{\mathrm{rep}} \cdot M \cdot d_s,
\label{eq:sea-fuel-route}
\end{equation}
em que \(\mathcal{S}^{*}_{o,d}\) é a cadeia selecionada, \(I_{o,d}^{\mathrm{rep}}\) é a mediana ponderada de todas as viagens elegíveis do par e \(d_s\) é a distância do subtrecho em milhas náuticas. A carga ANTAQ permanece como dado diagnóstico e como peso da intensidade; ela não substitui a carga informada no cenário. As intensidades resolvidas para o corredor selecionado também permanecem registradas para auditoria, mas não substituem o estimador representativo aplicado ao cenário.

O modelo segue a seguinte hierarquia:

\begin{enumerate}
  \item resolve a intensidade de cada viagem por IMO, classe ou tipo, nessa ordem;
  \item agrega todos os subtrechos de cada viagem e reúne todas as viagens elegíveis do mesmo par origem--destino;
  \item calcula a mediana das intensidades de viagem ponderada pelo trabalho de transporte;
  \item seleciona entre corredores completos observados sem costurar viagens distintas, usando essa escolha somente para caminho e distância;
  \item atribui combustível à carga do cenário aplicando a intensidade do par à soma das distâncias selecionadas;
  \item em usos específicos de decomposição, mantém submodelos de hoteling e operações portuárias.
\end{enumerate}

Uma consequência importante dessa hierarquia é a decisão de não somar hoteling separadamente quando a avaliação marítima já está baseada no KPI direcional por \emph{transport work}. Como hipótese de fronteira, o projeto trata esse indicador anual como abrangente do consumo específico reportado e evita somar hoteling novamente para reduzir o risco de dupla contagem.

\section{Benchmark com o workbook associado aos estudos de Gustavo Costa}

O benchmark explícito confronta o modelo com a planilha \texttt{Dados Relatorio 2.xlsx}, enviada como material de referência e usada no projeto como artefato de apoio às comparações e aos parâmetros operacionais \citep{workbookdados}. O pipeline:

\begin{enumerate}
  \item lê e normaliza os pares origem--destino presentes no workbook;
  \item converte o benchmark para a unidade operacional de \num{1} contêiner e \num{14} toneladas;
  \item confronta, quando possível, os mesmos pares com o motor do CabotageLens.
\end{enumerate}

O parser identifica \num{21} pares comparáveis entre \num{6} cidades. O benchmark documental produz duas referências relevantes: emissões semanais totais do workbook de \num{7614.97} tCO2e para o cenário rodoviário e \num{4159.79} tCO2e para o cenário de cabotagem, o que representa redução agregada de aproximadamente \num{45.4}\%. Quando se observa a média simples da economia percentual por par, o valor é de \num{46.7}\%. Um exemplo emblemático é Manaus--Fortaleza, para o qual o workbook indica \num{1733.9} kgCO2e por contêiner na rota rodoviária e \num{751.6} kgCO2e por contêiner na alternativa de cabotagem, isto é, uma redução de \num{56.7}\%.

Esse benchmark não deve ser lido como validação final exaustiva do projeto, mas como comparação de ordem de grandeza entre uma referência documental e o motor do repositório. Seu valor metodológico está no caminho auditável para comparar resultados publicados, artefatos de planilha e resultados do aplicativo em uma mesma unidade funcional.

\chapter{Resultados e discussão}
\label{chap:resultados}

\section{Evidências empíricas do modelo}

As evidências empíricas do CabotageLens abrangem a atividade observada, a intensidade energética e a matriz direcional.

A base processada da ANTAQ representa viagens, paradas e cargas conteinerizadas em 2025. Para cada operação elegível, ela identifica o navio, a sequência de portos e os movimentos de carga em cada parada.

O modelo marítimo associa indicadores anuais reportados pelo EU MRV por IMO e utiliza um KPI de combustível por \emph{transport work}. Quando o IMO não possui correspondência positiva, a intensidade recebe um fallback estatístico de classe ou tipo. A escolha do KPI evita confundir consumo absoluto anual do navio com intensidade normalizada pelo serviço de transporte.

A \texttt{sea\_matrix} reúne distâncias, intensidade energética direcional, alternativas de corredor e proveniência. Para cada par origem--destino, o motor dispõe de uma mediana das intensidades de viagem ponderada pelo trabalho de transporte, enquanto os corredores observados permanecem como alternativas auditáveis de caminho e distância.

\section{Interpretação da cobertura ANTAQ+MRV}

Uma cobertura de \num{62.5}\% dos IMOs e \num{59.5}\% das viagens significa que a evidência específica por navio permanece parcial. Ela cobre mais da metade da carga observada tanto em toneladas quanto em TEU, mas não autoriza tratar todas as intensidades como medições dos navios presentes na ANTAQ.

As \num{536} viagens sem correspondência positiva por IMO recebem um fallback explicitamente identificado de tipo; o fallback por classe também é suportado. Por isso, a resolução numérica alcança \num{100.0}\%, enquanto a cobertura exata continua parcial. A interpretação deve preservar essa diferença: o sistema preenche a lacuna numérica por uma regra estatística rastreável, sem apresentar o fallback como observação individual do navio ausente.

\section{Interpretação do benchmark com Gustavo Costa}

Os estudos associados a Gustavo Costa são particularmente úteis como referência para este trabalho porque combinam duas coisas que também orientam o CabotageLens: visão de rede multimodal e preocupação com desempenho energético e ambiental da cabotagem brasileira \citep{competitiveness2024,decarb2024}. A planilha local associada a esse conjunto de trabalhos, tratada no projeto como material de benchmark operacional, reforça a ideia de que corredores longos podem produzir vantagens expressivas para a cabotagem, especialmente quando a unidade funcional é o contêiner e quando a comparação é feita em termos de emissões diretas por remessa.

O fato de o workbook indicar economia média da ordem de \num{46.7}\% por par comparável é consistente com a literatura de \emph{short sea shipping} e com a expectativa de que corredores longos e costeiros favoreçam o transporte marítimo \citep{shortsea2019}. O benchmark exige a verificação da ordem de grandeza produzida pelo motor diante das referências documentais. O script do repositório executa essa comparação de forma reprodutível.

\section{Contribuição metodológica}

A contribuição metodológica do CabotageLens é um modelo marítimo baseado em atividade de viagens observadas, intensidade específica por IMO quando disponível e fallback estatístico registrado quando necessário. Os modelos por classe ou tipo cumprem a função de fallback, enquanto a fonte primária é a associação individual por IMO.

Do ponto de vista acadêmico, a formulação é compatível com a lógica dos artigos de referência usados neste trabalho. O estudo de competitividade em super-rede sugere que a unidade analítica correta é a rede multimodal observável \citep{competitiveness2024}. O estudo de descarbonização, por sua vez, enfatiza a necessidade de rigor no tratamento dos combustíveis e da unidade funcional \citep{decarb2024}. O cruzamento ANTAQ+MRV operacionaliza essas duas exigências em um motor reproduzível.

\chapter{Implementação computacional}
\label{chap:implementacao}

\section{Camadas do sistema}

O projeto está organizado em três camadas principais. A primeira é a camada de domínio, concentrada em \texttt{modules/}, responsável por roteamento, avaliação rodoviária, avaliação marítima, custos, emissões, persistência e serviços auxiliares. A segunda é a camada de orquestração, composta pelos \texttt{scripts/} de execução reproduzível. A terceira é a camada de interface, composta pela aplicação Streamlit.

Essa separação permite que o mesmo motor alimente diferentes modos de uso: comparação unitária, avaliação em lote, dashboard de cabotagem e benchmark com planilhas de referência. Também reduz o risco de divergência metodológica entre o que o usuário vê na interface e o que o pipeline realmente calcula.

\section{Persistência, banco e artefatos}

O projeto utiliza Supabase Postgres como backend persistente e Supabase Storage como repositório de ativos de dados e logs quando configurado. As rotas rodoviárias podem ser reutilizadas a partir do cache em Postgres, reduzindo chamadas redundantes à API de roteamento. As tabelas derivadas da ANTAQ foram materializadas tanto em CSV quanto em estrutura relacional (\texttt{antaq\_voyages}, \texttt{antaq\_voyage\_stops}, \texttt{antaq\_voyage\_stop\_calls} e \texttt{antaq\_voyages\_raw}), o que facilita tanto a auditoria quanto o consumo por dashboards.

\section{Interface e reprodutibilidade}

A aplicação expõe, entre outras funcionalidades, uma avaliação unitária, avaliações em lote, visualizações geográficas, \emph{heatmap} e um dashboard específico de cabotagem. Os componentes visuais leem a mesma base observada e os mesmos artefatos processados usados pelo motor de cálculo, reforçando a coerência entre análise, visualização e documentação.

\chapter{Limitações e ameaças à validade}
\label{chap:limitacoes}

O escopo do projeto possui limitações que condicionam a interpretação dos resultados.

Em primeiro lugar, a base observada está concentrada no recorte de 2025. Isso implica que a amostra de viagens, ainda que grande, não representa automaticamente a dinâmica de outros anos. Além disso, como toda janela de observação temporal limitada, a segmentação por IMO pode produzir primeiras e últimas viagens parciais.

Em segundo lugar, a cobertura exata por IMO no MRV é parcial. Das \num{1324} viagens, \num{788} usam intensidade específica por IMO e \num{536} dependem da média aparada em 1\% por cauda do tipo \emph{container ship}; o fallback por classe também é suportado, mas não é acionado pelos metadados disponíveis nesse recorte. Para o tipo \emph{container ship}, a regra parte de \num{243} valores, exclui dois em cada cauda e retém \num{239}, resultando em \num{9.322050} g/(t$\cdot$nm), diante de uma média aritmética bruta de \num{21.661852} e mediana bruta de \num{4.620000}. A regra de 1\% é uma política geral de tratamento de outliers e não é calibrada para reproduzir um corredor específico. A taxa de resolução numérica é, portanto, completa, mas a qualidade da evidência não é homogênea. A aparagem reduz a influência de valores extremos sem transformar o fallback em observação do navio ausente; resultados com fallback continuam exigindo incerteza maior do que aqueles associados diretamente ao IMO.

No par Santos--Manaus, as \num{89} observações viagem--OD de \num{22} corredores possuem trabalho positivo. A mediana ponderada de todas elas é \num{9.322050} g/(t$\cdot$nm): \num{40} usam IMO exato e \num{49} usam o fallback de tipo, que concentra \num{59.54}\% do trabalho de transporte. Essas 49 observações idênticas atravessam o percentil ponderado de 50\% e, portanto, controlam o escalar final. O escopo computacional cobre todas as viagens elegíveis, mas isso não equivale a 89 intensidades de combustível medidas de forma independente. A média aritmética ponderada das mesmas observações é \num{25.243618} g/(t$\cdot$nm) e aparece apenas como diagnóstico da influência da cauda superior.

Em terceiro lugar, a escolha de porto por proximidade geográfica ainda é uma heurística simplificadora. Em termos reais, a escolha do porto depende de oferta de serviço, contrato logístico, tempo de trânsito, custo portuário, disponibilidade de equipamento e frequência de escalas. Incorporar esses fatores exigirá uma camada adicional de modelagem de serviço, não apenas de geometria.

Em quarto lugar, o escopo ambiental é predominantemente \emph{tank-to-wheel}. Essa fronteira atende à comparação operacional de combustível e emissões diretas, mas não substitui análises \emph{well-to-wake} ou de ciclo de vida completo para cenários com HVO, LNG, biodiesel ou eletrificação portuária \citep{lifecycle2020,maritimelca2024,decarb2024}.

Por fim, o benchmark com o workbook de referência é um instrumento de comparação direcional e de ordem de grandeza, não um veredito definitivo. A equivalência exata entre workbook, artigo publicado e motor do aplicativo depende de comparação pareada sob as mesmas hipóteses de serviço, distância, carga, alocação e fronteira ambiental.

\chapter{Conclusão}
\label{chap:conclusao}

O CabotageLens combina quatro elementos em uma ferramenta de análise aplicada: um motor rodoviário explícito e auditável; uma base observada de cabotagem brasileira construída a partir da ANTAQ; um mecanismo de associação por IMO com indicadores empíricos do EU MRV; e uma matriz marítima com mediana direcional ponderada por trabalho de transporte para cada par origem--destino.

O resultado substantivo é uma infraestrutura analítica capaz de vincular cada número às evidências que o sustentam. A intensidade de cada par é estimada pela mediana ponderada das viagens que efetivamente ligaram sua origem ao destino na base observada; a avaliação unitária utiliza uma unidade explícita de \num{1} contêiner e \num{14} toneladas; e o workbook de referência é comparado ao motor por um roteiro programático e reproduzível.

As extensões prioritárias são ampliar o recorte temporal da base observada, aumentar a cobertura MRV, incorporar oferta real de serviço marítimo, expandir a comparação com benchmarks publicados e desenvolver análises \emph{well-to-wake}. Dentro da fronteira documentada, o CabotageLens oferece uma base auditável para análises comparativas de rodoviário versus cabotagem no Brasil.

\bibliographystyle{unsrtnat}
\IfFileExists{docs/references.bib}{\bibliography{docs/references}}{\bibliography{references}}

\end{document}
